\documentclass[11pt]{article}
\usepackage{amsmath,amssymb,amsthm}
\usepackage[margin=1in]{geometry}

\newtheorem{theorem}{Theorem}
\newtheorem{remark}{Remark}

\begin{document}

\title{Unique Finite-Radius Exterior Continuation of the\\
Exceptional $\ell=2$ Axial GfE Solution}
\date{7 August 2026}
\maketitle

\paragraph{Scope.}
This theorem builds on the certified convergent exceptional-horizon
logarithmic Frobenius solution at
\[
  \omega=\frac{i}{4M}.
\]
It proves unique continuation through every finite radius in the open
Schwarzschild exterior.  It does not classify the singular endpoint
$r=\infty$, impose an infinity boundary condition, establish a global
exceptional mode, or imply physical instability.

\begin{theorem}[No finite-radius obstruction and unique exterior continuation]
Fix $M>0$ and $\beta\neq0$, and set
\[
  \omega=\frac{i}{4M}.
\]
For the exact uneliminated $\ell=2$ axial GfE equations, the
Douglis--Nirenberg derivative staircase is
\[
  E_0:\ (H_0^{(4)},H_1^{(3)}),
  \qquad
  E_1:\ (H_0^{(3)},H_1^{(2)}).
\]
Its exact principal matrix is
\[
C(r)=
\begin{pmatrix}
-\dfrac{8M\beta^2(r-2M)}{r^4}
&
 \dfrac{2\beta^2(r-2M)}{r^4}
\\[7pt]
 \dfrac{2\beta^2(r-2M)}{r^4}
&
-\dfrac{\beta^2(r-2M)
(256M^3-128M^2r+r^3)}
{2Mr^7}
\end{pmatrix},
\]
with determinant
\[
  \det C(r)
  =
  -\frac{512M^2\beta^4(r-2M)^3}{r^{11}}.
\]
Hence
\[
  \det C(r)\neq0
  \qquad
  \text{for every finite }r>2M.
\]

More strongly, the equations admit an explicit first-order reduction
on the whole open exterior that never divides by a coefficient having
an exterior zero.  The coefficient of $H_0^{(3)}$ in $E_1$ is
\[
  a_3(r)
  =
  \frac{2\beta^2(r-2M)}{r^4},
\]
so $E_1=0$ uniquely solves for $H_0^{(3)}$ at every $r>2M$.
Differentiating this constraint and combining it with $E_0=0$ gives
the remaining pivot
\[
  \kappa(r)
  =
  \frac{256M^2\beta^2(r-2M)^2}{r^7},
\]
which is also nonzero for every $r>2M$ and uniquely solves for
$H_1^{(3)}$.

Therefore, with
\[
  Z=
  (H_0,H_0',H_0'',H_1,H_1',H_1'')^T,
\]
the exact uneliminated system is equivalent on $r>2M$ to a linear
first-order system
\[
  Z'(r)=\mathcal A(r)Z(r),
\]
whose coefficients are rational functions with no finite poles in the
open exterior.  In fact the symbolic coefficient census shows that
all finite coefficient poles lie in
\[
  \{0,2M\}.
\]

Consequently, for every
\[
  2M<r_0<R<\infty
\]
and every admissible trace $Z(r_0)$, standard linear ODE existence and
uniqueness gives a unique solution on $[r_0,R]$.  Since $R$ is
arbitrary, every local solution at $r_0>2M$ extends uniquely through
every finite exterior radius.

In particular, the previously certified convergent exceptional
horizon solution extends uniquely from its punctured horizon
neighborhood to the entire open exterior
\[
  2M<r<\infty .
\]
\end{theorem}

\paragraph{Proof.}
The displayed staircase coefficients, determinant, safe pivot
$a_3$, and differentiated-constraint pivot $\kappa$ are exact symbolic
consequences of the current canonical GfE source at
$\omega=i/(4M)$.  The denominator census checks every jet coefficient
appearing in $E_0$ and $E_1$, as well as the two inverse pivots used in
the reduction, and finds no finite denominator zero for $r>2M$.
Thus $\mathcal A(r)$ is continuous, indeed analytic, on every compact
interval $[r_0,R]\subset(2M,\infty)$.

A linear first-order system with continuous coefficients has a unique
solution on each such compact interval and cannot develop a
finite-radius blow-up obstructing continuation.  Uniqueness makes the
solutions obtained on nested compact intervals agree on overlaps.
Taking the union over all finite $R$ proves unique continuation on the
whole open exterior.

The convergent Frobenius theorem supplies a genuine solution on some
interval $(2M,2M+\rho)$.  Choosing any
$r_0\in(2M,2M+\rho)$ and applying the preceding continuation result
extends that same solution uniquely to every finite $r>2M$.
\qed

\begin{remark}[What remains]
The endpoint $r=\infty$ is not a finite-radius degeneracy and is not
covered by compact-interval continuation.  The next unresolved
problem is therefore to derive the exact infinity asymptotic basis and
decompose the uniquely continued exceptional solution into the
allowed and forbidden infinity channels.
\end{remark}

\begin{remark}[No instability claim]
Unique exterior continuation does not by itself imply that the
exceptional solution satisfies an outgoing, decaying, finite-energy,
or quasinormal boundary condition at infinity.  No global exceptional
mode or physical instability is asserted here.
\end{remark}

\end{document}
